\begin{lemma}
\label{lemma-characterize-geometrically-integral}
Let $k$ be a field. Let $S$ be a $k$-algebra.
The following are equivalent
\begin{enumerate}
\item $S$ is geometrically integral over $k$,
\item for every finite extension $k'/k$ of fields
the ring $S \otimes_k k'$ is a domain,
\item $S \otimes_k \overline{k}$ is a domain
where $\overline{k}$ is the algebraic closure of $k$.
\end{enumerate}
\end{lemma}

\begin{proof}
Condition (1) implies (2) by
restricting the field extensions.
Assume (2). In particular its
identity extension makes $S$
a nonzero domain. Choose the
original algebraic closure
$\overline k$ and consider its
finite subextensions $E/k$.
They are directed under compositum
and contain every finite list
of coefficients of a tensor.
The original injective coefficient
maps therefore give
$$
S\otimes_k\overline k=\bigcup_E S\otimes_k E.
$$
Each ring in this union is a
domain by (2). The union is
nonzero, since the maps preserve
the nonzero unit. For two
nonzero elements, choose one
stage containing both. Their
product there is nonzero and
remains so by injectivity.
This proves (3), including
nonemptiness and all zero-divisor
conditions.

Assume (3), and let $L/k$
be any original extension.
Choose a prime $\mathfrak a$
of the nonzero field tensor
$L\otimes_k\overline k$ and
form the compatible receiving
field
$$
F=\operatorname{Frac}((L\otimes_k\overline k)/\mathfrak a).
$$
The maps $a\mapsto[a\otimes1]$
and $b\mapsto[1\otimes b]$
embed the two original fields
and agree on $k$. Write
$B=S\otimes_k\overline k$.
This is a nonzero domain
by (3). Since $\overline k$
is perfect, Lemma
\ref{lemma-perfect-reduced}
makes $B\otimes_{\overline k}F$
reduced. Since $\overline k$
is separably closed and both
$B$ and $F$ have irreducible
spectra, Lemma
\ref{lemma-separably-closed-irreducible}
makes its spectrum irreducible.
Thus it is a domain by Lemma
\ref{lemma-geometrically-integral}'s
underlying ring criterion.

For either original subfield
$H=L$ or $H=\overline k$ of
the receiving $F$, the maps
$$
(S\otimes_k H)\otimes_H F\longrightarrow S\otimes_k F,
\qquad(s\otimes a)\otimes b\longmapsto s\otimes ab
$$
and $s\otimes b\mapsto(s\otimes1)\otimes b$
are inverse by balancing.
They identify the domain just
proved with $S\otimes_k F$.
The original scalar map
$S\otimes_k L\to(S\otimes_k L)\otimes_L F$
is injective by extension of
an $L$-vector-space basis.
Its source is therefore a
nonzero subring of that domain,
so is a domain. This proves
(1) for every original $L/k$.

The same receiving-field maps
prove geometric integrality
is preserved and reflected
by every field extension
$L/k$. Preservation follows
from the displayed comparison
for all $F/L$. For reflection,
suppose $S\otimes_k L$ is
geometrically integral over $L$.
For any $E/k$, construct a
compatible common field for
$E$ and $L$ as above. The
hypothesis makes $S\otimes_k F$
a domain, and the injection
$S\otimes_k E\to S\otimes_k F$
then makes the original
$S\otimes_k E$ a domain.
This is precisely geometric
integrality of $S/k$.
The zero algebra satisfies
none of (1)--(3), as each
condition includes a nonzero
domain after scalar extension.
\end{proof}